%% Question bank chapter file - generated by bank_tools.py build (do not edit while the bank is still being built:
%% it is regenerated from the per-test data in Question Bank\data). Every question is one bankq block:
%%   \begin{bankq}[id=..., source={...}, paper=.., qno=.., type=..., subtopic={...}, answer={...}, answer_src=key|solved,
%%                 working={...}, difficulty=1-10, diff_parts={K S I C F}, diff_note={...},
%%                 group=<passage id or empty>, also_in={duplicate copies in other tests}]
%% difficulty = min(10, K+S+I+C+F-1): K concepts 1-3, S steps 1-4, I insight 0-2, C computation 0-2, F format 0-1
%%   ...question LaTeX, options, figures (\bfig / \bfigi -> img/)...
%%   \end{bankq}
%% Passages shared by several questions are bankpassage blocks; their questions carry group=<passage id>.
\documentclass[10pt]{article}
\usepackage[a4paper,margin=1.8cm]{geometry}
\usepackage[T1]{fontenc}
\usepackage{amsmath,amssymb,mathtools}
\usepackage{graphicx}
\usepackage[export]{adjustbox}
\usepackage{xcolor,needspace,keyval}
\usepackage{enumitem}
\usepackage{booktabs,array,tabularx}
\usepackage[hidelinks]{hyperref}
\setlength{\parindent}{0pt}
\setlength{\parskip}{3pt}
\makeatletter
\define@key{bq}{id}{\def\bq@id{#1}}\define@key{bq}{source}{\def\bq@src{#1}}\define@key{bq}{paper}{\def\bq@paper{#1}}
\define@key{bq}{qno}{\def\bq@qno{#1}}\define@key{bq}{type}{\def\bq@type{#1}}\define@key{bq}{subtopic}{\def\bq@sub{#1}}
\define@key{bq}{answer}{\def\bq@ans{#1}}\define@key{bq}{answer_src}{\def\bq@asrc{#1}}\define@key{bq}{working}{\def\bq@work{#1}}
\define@key{bq}{group}{\def\bq@group{#1}}\define@key{bq}{also_in}{\def\bq@also{#1}}
\define@key{bq}{difficulty}{\def\bq@diff{#1}}\define@key{bq}{diff_parts}{\def\bq@dparts{#1}}\define@key{bq}{diff_note}{\def\bq@dnote{#1}}
\newenvironment{bankq}[1][]{\def\bq@work{}\def\bq@also{}\def\bq@group{}\def\bq@diff{?}\def\bq@dparts{}\def\bq@dnote{}\setkeys{bq}{#1}%
  \par\needspace{4\baselineskip}\medskip\noindent{\small\textbf{\bq@id}\hfill\color{gray}[\bq@src, Paper \bq@paper, Q\bq@qno]
  \bq@type\ $\cdot$ \bq@sub\ $\cdot$ \textbf{D\bq@diff}}\par\nopagebreak\relax}
  {\par{\small\color{teal}Answer: \textbf{\bq@ans}\ifx\bq@work\@empty\else\ (solved: \bq@work)\fi
  \ifx\bq@also\@empty\else\par\color{gray}Also in: \bq@also\fi}\par}
\define@key{bp}{id}{\def\bp@id{#1}}\define@key{bp}{source}{\def\bp@src{#1}}\define@key{bp}{questions}{\def\bp@qs{#1}}
\newenvironment{bankpassage}[1][]{\setkeys{bp}{#1}\par\needspace{4\baselineskip}\bigskip\noindent
  {\small\textbf{Passage \bp@id}\hfill\color{gray}[\bp@src; questions \bp@qs]}\par\nopagebreak\relax\itshape}{\par}
\makeatother
\newcommand{\bfig}[2]{\par\smallskip{\centering\includegraphics[width=#1pt,max width=0.62\linewidth,max height=0.32\textheight]{img/#2}\par}\smallskip}
\newcommand{\bfigi}[2]{\includegraphics[width=#1pt,max width=\linewidth,max height=0.32\textheight,valign=c]{img/#2}}
\newcommand{\bmissing}[1]{\textit{[figure missing in the source conversion: \detokenize{#1}]}}
\newenvironment{qoptions}{\begin{itemize}[leftmargin=1.8em,itemsep=1pt,topsep=2pt,label={}]}{\end{itemize}}
\newcommand{\opt}[2]{\item[(#1)] #2}
\providecommand{\fwstart}{}\providecommand{\fwend}{}\providecommand{\qfullsplit}{}
\newcolumntype{L}[1]{>{\raggedright\arraybackslash}p{#1}}
\newcolumntype{C}[1]{>{\centering\arraybackslash}p{#1}}
\newcolumntype{R}[1]{>{\raggedleft\arraybackslash}p{#1}}
\providecommand{\sectionhdr}[1]{}\providecommand{\subsectionhdr}[1]{}


\begin{document}
{\Large\bfseries Complex Numbers --- Question Bank}\par\smallskip
{\small Mathematics chapter M02. Questions from all converted test papers, grouped by source series. Answers: printed key where it matches the paper, otherwise solved (answer\_src=solved, with working). Exact repeats across papers appear once, with also\_in listing the other copies.}\par\medskip\hrule

\section*{Score-I papers}
\begin{bankq}[id=SI12-P1-MATH-S1-Q01, source={Score-I - Test 12 (09 Mar 2022)}, paper=1, qno=1, type=multi, subtopic={De Moivre}, answer={A,D}, answer_src=key, difficulty=5, diff_parts={K2 S2 I1 C0 F1}, diff_note={concepts: roots of a quadratic as complex numbers, de Moivre's theorem | steps: roots 1+-i and x + alpha = (cos theta + i sin theta)/sin theta; raise to n and subtract; divide by 2i and match the options | I: recognise the polar form from cot theta}]
If $\alpha,\beta$ be the roots of the equation $\mu^2-2\mu+2=0$ and if $\cot\theta=x+1$, then $\dfrac{(x+\alpha)^n-(x+\beta)^n}{\alpha-\beta}$ is equal to, {(where $\theta\ne k_1\pi$, $k_1\in I$)}
\begin{qoptions}
\opt{A}{$\dfrac{\sin n\theta}{\sin^n\theta}$}
\opt{B}{$\dfrac{\cos n\theta}{\cos^n\theta}$}
\opt{C}{$\dfrac{\sin n\theta}{\cos^n\theta}$}
\opt{D}{$\dfrac{\operatorname{cosec}^n\theta}{\operatorname{cosec}n\theta}$, where $n\theta\ne m\pi$, $m\in I$}
\end{qoptions}
\end{bankq}
\begin{bankq}[id=SI12-P2-MATH-S2-Q04, source={Score-I - Test 12 (09 Mar 2022)}, paper=2, qno=4, type=single, subtopic={Geometry}, answer={12.00}, answer_src=key, difficulty=5, diff_parts={K2 S3 I1 C0 F0}, diff_note={concepts: argument curves as rays, tangent circle | steps: write the four rays as lines; distance from (0,alpha) to the line equals 4 sqrt 2; alpha = 6 or -10; check which tangent points lie on the rays; sum  | alpha}]
Consider curves $C_1:|\arg(z-2)|=\dfrac\pi4$; $C_2:|\arg(z+2)|=\dfrac{3\pi}4$ and $C_3:|z-\alpha i|=4\sqrt2$, $\alpha\in R$ on the complex plane such that $C_3$ touches both $C_1$ and $C_2$, then sum of absolute value(s) of $\alpha$ is
\end{bankq}

\section*{Nurture SRG papers}
\begin{bankq}[id=NSRG03-P1-MATH-S1-Q01, source={Nurture SRG - Test 03 (04 Feb 2024)}, paper=1, qno=1, type=single, subtopic={Roots of unity}, answer={B}, answer_src=solved, working={$\sum\dfrac1{\omega^r-1}=-\dfrac{n-1}2$ with $n=51$, so $-25$}, difficulty=6, diff_parts={K1 S2 I2 C2 F0}, diff_note={concepts: sum of reciprocals over roots of unity | steps: roots of the equation; the identity for sums over roots of unity; evaluate}]
If $Z_r$; $r=1,2,3,\dots,50$ are the roots of the equation $\displaystyle\sum_{r=0}^{50}(Z)^r=0$, then the value of $\displaystyle\sum_{r=1}^{50}\dfrac1{Z_r-1}$ is
\begin{qoptions}
\opt{A}{$-85$}
\opt{B}{$-25$}
\opt{C}{$25$}
\opt{D}{$75$}
\end{qoptions}
\end{bankq}
\begin{bankq}[id=NSRG03-P1-MATH-S1-Q02, source={Nurture SRG - Test 03 (04 Feb 2024)}, paper=1, qno=2, type=single, subtopic={Modulus and argument}, answer={A}, answer_src=solved, working={Moduli are 1; $z^2=w^4$ and $z^3=-w^5$ give $z=-w$, $w^2=1$, so both real}, difficulty=6, diff_parts={K1 S2 I2 C2 F0}, diff_note={concepts: complex numbers satisfying two equations | steps: moduli from the equations; relation of the arguments; solve; nature of z and w}]
If $z$ \& $w$ are two complex numbers simultaneously satisfying the equations, $z^3+w^5=0$ and $z^2\overline w^4=1$, then:
\begin{qoptions}
\opt{A}{$z$ and $w$ both are purely real}
\opt{B}{$z$ is purely real and $w$ is purely imaginary}
\opt{C}{$w$ is purely real and $z$ is purely imaginary}
\opt{D}{$z$ and $w$ both are imaginary}
\end{qoptions}
\end{bankq}
\begin{bankq}[id=NSRG03-P1-MATH-S1-Q05, source={Nurture SRG - Test 03 (04 Feb 2024)}, paper=1, qno=5, type=single, subtopic={Geometry}, answer={C}, answer_src=solved, working={Numerical check of the expression for an acute triangle gives $-1$}, difficulty=6, diff_parts={K1 S2 I2 C2 F0}, diff_note={concepts: expression with the circumcentre of a triangle | steps: ratio of vectors from the circumcentre; angles; use of the sine of double angles; sum}]
If $z_1,z_2$ and $z_3$ are the vertices of $\Delta ABC$, which is not right angled triangle taken in anti-clock wise direction and $z_0$ is the circumcentre, then $\left(\dfrac{z_0-z_1}{z_0-z_2}\right)\dfrac{\sin2A}{\sin2B}+\left(\dfrac{z_0-z_3}{z_0-z_2}\right)\dfrac{\sin2C}{\sin2B}$ is equal to
\begin{qoptions}
\opt{A}{0}
\opt{B}{1}
\opt{C}{$-1$}
\opt{D}{2}
\end{qoptions}
\end{bankq}
\begin{bankq}[id=NSRG03-P1-MATH-S1-Q06, source={Nurture SRG - Test 03 (04 Feb 2024)}, paper=1, qno=6, type=single, subtopic={Geometry}, answer={C}, answer_src=solved, working={$OP=r$, $OQ=\sqrt2r$ with $45^\circ$ between them, so $PQ=r$ and the triangle is right isosceles}, difficulty=5, diff_parts={K1 S2 I2 C1 F0}, diff_note={concepts: triangle formed by two complex numbers and the origin | steps: lengths from the moduli; angle between them; cosine rule; classify}]
Let $P$ denotes a complex number $z=r(\cos\theta+i\sin\theta)$ on the Argand's plane, and $Q$ denotes a complex number $\sqrt{2|z|^2}\left(\cos\left(\theta+\dfrac\pi4\right)+i\sin\left(\theta+\dfrac\pi4\right)\right)$. If `O' is the origin, then $\Delta OPQ$ is
\begin{qoptions}
\opt{A}{Isosceles but not right angled}
\opt{B}{Right angled but not isosceles}
\opt{C}{Right isosceles}
\opt{D}{Equilateral}
\end{qoptions}
\end{bankq}
\begin{bankq}[id=NSRG03-P1-MATH-S2-Q07, source={Nurture SRG - Test 03 (04 Feb 2024)}, paper=1, qno=7, type=multi, subtopic={Modulus and argument}, answer={A}, answer_src=solved, working={$|z_1|=13$: maximum of $|z_1+iz_2|$ is 17; the minimum and maximum of the quotient are $\tfrac{13}5$ and $\tfrac{13}3$}, difficulty=7, diff_parts={K1 S2 I2 C2 F1}, diff_note={concepts: maximum and minimum moduli of complex expressions | steps: triangle inequality; extreme values of each; check each option}]
If $z_1=5+12i$ and $|z_2|=4$ then
\begin{qoptions}
\opt{A}{Maximum $\left(|z_1+iz_2|\right)=17$}
\opt{B}{Minimum $\left(|z_1+(1+i|z_2|\,|\right)=13-9\sqrt2$}
\opt{C}{Minimum $\left|\dfrac{z_1}{z_2+\frac4{z_2}}\right|=\dfrac{13}4$}
\opt{D}{Maximum $\left|\dfrac{z_1}{z_2+\frac4{z_2}}\right|=\dfrac{13}4$}
\end{qoptions}
\end{bankq}
\begin{bankq}[id=NSRG03-P1-MATH-S2-Q08, source={Nurture SRG - Test 03 (04 Feb 2024)}, paper=1, qno=8, type=multi, subtopic={Geometry}, answer={A,D}, answer_src=solved, working={With $\angle O=2\phi$, $\dfrac{Z_2^2}{Z_3Z_1}=\dfrac{\cos^2\phi}{\cos2\phi}=\tfrac32$ gives $\cos2\phi=\tfrac12$: angles $\tfrac\pi6,\tfrac\pi3,\tfrac\pi2$ and area $\tfrac{\sqrt3}2R^2$ (read as $Z_1$ in the denominator)}, difficulty=8, diff_parts={K1 S3 I2 C2 F1}, diff_note={concepts: right triangle with the angle bisector and concyclic points | steps: positions of the points on the circle; use of the bisector; ratio condition; angles; area}]
Line $OQ$ is angle bisector of angle $O$ of right angle triangle $OPR$, right angle at $P$. Point $Q$ is such that $ORQP$ is concyclic. If point $O$ is origin and points $P,Q,R$ are represented by the complex numbers $Z_3$, $Z_2$, $Z_1$ respectively. If $\dfrac{Z_2^2}{Z_3Z_2}=\dfrac32$ then ($R$ is circum radius of $\Delta OPR$)
\begin{qoptions}
\opt{A}{Angles of $\Delta OPR$ are $\dfrac\pi6,\dfrac\pi3,\dfrac\pi2$}
\opt{B}{Angles of $\Delta OPR$ are $\dfrac\pi4,\dfrac\pi4,\dfrac\pi2$}
\opt{C}{Area of $\Delta OPR$ is $2\sqrt2R^2$}
\opt{D}{Area of $\Delta OPR$ is $\dfrac{\sqrt3}2R^2$}
\end{qoptions}
\end{bankq}
\begin{bankq}[id=NSRG03-P1-MATH-S2-Q09, source={Nurture SRG - Test 03 (04 Feb 2024)}, paper=1, qno=9, type=multi, subtopic={Geometry}, answer={A,B}, answer_src=solved, working={$\angle A=90^\circ$ gives $\mathrm{Re}(W_1\overline{W_2})=|W_1|^2$ and $\mathrm{Re}(W_2/W_1)=1$}, difficulty=6, diff_parts={K1 S2 I2 C1 F1}, diff_note={concepts: right angle at a vertex in terms of complex numbers | steps: vectors of the sides; perpendicular condition; real part relations; check}]
Let $z_1,z_2,z_3$ be the coordinate of verticles A, B, C of the triangle if $W_1=z_1-z_2$, $W_2=z_3-z_2$ and $\angle A=90^\circ$
\begin{qoptions}
\opt{A}{$\mathrm{Re}\left(\dfrac{W_2}{W_1}\right)=1$}
\opt{B}{$\mathrm{Re}\left(W_1\overline{W_2}\right)=|W_1|^2$}
\opt{C}{$\mathrm{Re}\left(W_1\overline{W_2}\right)=|W_2|^2$}
\opt{D}{$\mathrm{Re}\left(\dfrac{W_1}{W_2}\right)=1$}
\end{qoptions}
\end{bankq}
\begin{bankq}[id=NSRG03-P1-MATH-S3-Q03, source={Nurture SRG - Test 03 (04 Feb 2024)}, paper=1, qno=3, type=numerical, subtopic={Geometry}, answer={30}, answer_src=solved, working={Minimum of the sum of squared distances to $0$, $3$, $6i$ at the centroid $1+2i$ is $5+8+17=30$}, difficulty=6, diff_parts={K1 S2 I2 C1 F1}, diff_note={concepts: minimum of a sum of squared distances | steps: centroid; minimum value; evaluate}]
Find the minimum value of the expression $E=|z|^2+|z-3|^2+|z-6i|^2$ (where $z=x+iy$, $x,y\in R$)
\end{bankq}
\begin{bankq}[id=NSRG03-P2-MATH-S1-Q01, source={Nurture SRG - Test 03 (04 Feb 2024)}, paper=2, qno=1, type=single, subtopic={Roots of unity}, answer={A}, answer_src=solved, working={Adding and subtracting with the conjugate equation gives $\sum\dfrac1{a^2-a+1}=1$}, difficulty=6, diff_parts={K1 S2 I2 C2 F0}, diff_note={concepts: sum over reciprocal terms with a cube root of unity | steps: take the conjugate equation; add and subtract; identity; result}]
If $\dfrac1{a+\omega}+\dfrac1{b+\omega}+\dfrac1{c+\omega}+\dfrac1{d+\omega}=\dfrac1\omega$, where $a,b,c,d\in R$ and $\omega$ is a cube root of unity then $\sum\dfrac1{a^2-a+1}$ is equal to
\begin{qoptions}
\opt{A}{1}
\opt{B}{2}
\opt{C}{3}
\opt{D}{$\dfrac13$}
\end{qoptions}
\end{bankq}
\begin{bankq}[id=NSRG03-P2-MATH-S2-Q07, source={Nurture SRG - Test 03 (04 Feb 2024)}, paper=2, qno=7, type=multi, subtopic={Geometry}, answer={B,C,D}, answer_src=solved, working={$|z_1z_2+z_2z_3+z_3z_1|=|\bar z_1+\bar z_2+\bar z_3|$; the product expression is real; the isosceles right triangle gives zero real part}, difficulty=7, diff_parts={K1 S2 I2 C2 F1}, diff_note={concepts: statements on complex numbers on the unit circle | steps: angle subtended; identities with conjugates; reality of the expression; right isosceles triangle}]
Let $z_1,z_2$ and $z_3$ be three distinct complex numbers, satisfying $|z_1|=|z_2|=|z_3|=1$. Which of the following is/are true :
\begin{qoptions}
\opt{A}{If $\arg\left(\dfrac{z_1}{z_2}\right)=\dfrac\pi2$ then $\arg\left(\dfrac{z-z_1}{z-z_2}\right)>\dfrac\pi4$ where $|z|>1$}
\opt{B}{$|z_1z_2+z_2z_3+z_3z_1|=|z_1+z_2+z_3|$}
\opt{C}{$\mathrm{Im}\left(\dfrac{(z_1+z_2)(z_2+z_3)(z_3+z_1)}{z_1\cdot z_2\cdot z_3}\right)=0$}
\opt{D}{If $|z_1-z_2|=\sqrt2|z_1-z_3|=\sqrt2|z_2-z_3|$, then $\mathrm{Re}\left(\dfrac{z_3-z_1}{z_3-z_2}\right)=0$}
\end{qoptions}
\end{bankq}
\begin{bankq}[id=NSRG03-P2-MATH-S2-Q11, source={Nurture SRG - Test 03 (04 Feb 2024)}, paper=2, qno=11, type=multi, subtopic={Modulus and argument}, answer={A,B,C}, answer_src=solved, working={$\mathrm{Re}(z_1\bar z_2)=0$ makes $(c,d)=\pm(-b,a)$, so $|z_3|=|z_4|=1$ and $\mathrm{Re}(z_3\bar z_4)=0$}, difficulty=6, diff_parts={K1 S2 I2 C1 F1}, diff_note={concepts: complex numbers with perpendicular position vectors | steps: orthogonality in components; relation between a b c d; moduli; real and imaginary parts}]
If $z_1=a+ib$ and $z_2=c+id$ are complex number such that $|z_1|=|z_2|=1$ and $\mathrm{Re}\left(z_1\overline{z_2}\right)=0$, then the pair of complex numbers $z_3=a+ic$ and $z_4=b+id$ satisfies :
\begin{qoptions}
\opt{A}{$|z_3|=1$}
\opt{B}{$|z_4|=1$}
\opt{C}{$\mathrm{Re}\left(z_3\overline{z_4}\right)=0$}
\opt{D}{$\mathrm{Im}\left(z_3\overline{z_4}\right)=0$}
\end{qoptions}
\end{bankq}

\section*{TOAS papers}
\begin{bankq}[id=TOAS18-P1-MATH-S1-Q01, source={TOAS - Test 18 (30 Oct 2022)}, paper=1, qno=1, type=multi, subtopic={Geometry}, answer={A,B,C}, answer_src=solved, working={The line $x+y=\sqrt2$ meets the circle at $Z\approx(-0.76,2.18)$ with $\mathrm{Im}\ge1$, so $|Z|\approx2.3<3$; $|\omega|<3+\sqrt5\approx5.24$ so $|\omega|<4$ is possible and $|Z|-|\omega|>-6$}, difficulty=6, diff_parts={K1 S3 I1 C1 F1}, diff_note={concepts: geometry of loci in the complex plane | steps: locus line and circle; intersection point; modulus of Z; range of omega; compare}]
Let $A$, $B$, $C$ be the three sets of complex numbers as defined below $A=\{Z:\operatorname{Im}Z\ge1\}$, $B=\{Z:|Z-2-i|=3\}$ and $C=\left\{Z:\operatorname{Re}\big((1-i)Z\big)=\sqrt2\right\}$. Let `$Z$' be any point in $A\cap B\cap C$ and let `$\omega$' be any point satisfying $|\omega-2-i|<3$, then which of the following may be correct ?
\begin{qoptions}
\opt{A}{$|Z|<3$}
\opt{B}{$|\omega|<4$}
\opt{C}{$|Z|-|\omega|>-6$}
\opt{D}{$|Z|>3$}
\end{qoptions}
\end{bankq}
\begin{bankq}[id=TOAS18-P1-MATH-S3-Q03, source={TOAS - Test 18 (30 Oct 2022)}, paper=1, qno=3, type=numerical, subtopic={Algebra}, answer={6}, answer_src=solved, working={The ratio of the two column entries has modulus 1, so $1+(a_1^2-2b_1^2)\omega+([x]+[y]+[z])\omega^2=0$ forces $a_1^2-2b_1^2=1$ ($a_1=3$, $b_1=2$) and $[x]+[y]+[z]=1$; the required sum is $1+3+2=6$}, difficulty=6, diff_parts={K2 S2 I2 C0 F1}, diff_note={concepts: cube roots, cyclic combinations | steps: relation between the two combinations; coefficient comparison; prime solution; evaluate the greatest integer sum}]
$a,b,c$ are cube roots of $p$($p<0$) then for any permissible real value of $x,y,z$ it is given that $\begin{vmatrix}xa+yb+zc\\xb+yc+za\end{vmatrix}+\left(a_1^2-2b_1^2\right)\omega+([x]+[y]+[z])\omega^2=0$ where $\omega$ is the imaginary cube root of unity and $a_1,b_1$ are prime numbers. Then $[x+a_1]+[y+b_1]+[z]$ is equal to (where $[.]$ denotes G.I.F.)
\end{bankq}
\begin{bankq}[id=TOAS18-P1-MATH-S3-Q06, source={TOAS - Test 18 (30 Oct 2022)}, paper=1, qno=6, type=numerical, subtopic={De Moivre}, answer={10}, answer_src=solved, working={$Z=e^{-i\pi/6}$: $Z^{101}=e^{-5i\pi/6}$, $+i$ gives $e^{5i\pi/6}$; raised to 106: $e^{i\pi/3}$; need $-n\pi/6\equiv\pi/3\pmod{2\pi}$, so $n=10$}, difficulty=5, diff_parts={K1 S2 I1 C1 F1}, diff_note={concepts: powers of a complex number | steps: argument of Z; reduce the powers; sum; equation for n}]
If $Z=\dfrac12\left(\sqrt3-i\right)$ then the least positive integral value of `$n$' such that $\left(Z^{101}+i^{109}\right)^{106}=Z^n$ is (where $i^2=-1$)
\end{bankq}
\begin{bankq}[id=TOAS18-P2-MATH-S2-Q07, source={TOAS - Test 18 (30 Oct 2022)}, paper=2, qno=7, type=multi, subtopic={Geometry}, answer={B,C}, answer_src=key, difficulty=3, diff_parts={K1 S2 I0 C0 F1}, diff_note={concepts: loci in the complex plane | steps: disc and ellipse; distances between points; range}]
If $z_1$ is a point on locus given by $|z+8|\le2$ and $z_2$ is a point on locus given by $|z+3|+|z-3|\le10$ then $|z_1-z_2|$ can be
\begin{qoptions}
\opt{A}{0}
\opt{B}{2}
\opt{C}{10}
\opt{D}{18}
\end{qoptions}
\end{bankq}
\begin{bankq}[id=TOAS18-P2-MATH-S3-Q04, source={TOAS - Test 18 (30 Oct 2022)}, paper=2, qno=4, type=numerical, subtopic={Modulus and argument}, answer={4}, answer_src=key, difficulty=4, diff_parts={K1 S2 I1 C0 F1}, diff_note={concepts: real ratio of complex expressions | steps: imaginary part zero; condition on the modulus; evaluate}]
Let $z$ be a complex number with non-zero imaginary part. If $\dfrac{2+3z+4z^2}{2-3z+4z^2}$ is a real number, then the value of $8|z|^2$ is \underline{\hspace{1.2cm}}.
\end{bankq}

\section*{TMAS papers}
\begin{bankq}[id=TMAS04-P2-MATH-S1-Q01, source={TMAS - Test 04 (19 Nov 2023)}, paper=2, qno=1, type=multi, subtopic={Roots of unity}, answer={A,C}, answer_src=key, difficulty=6, diff_parts={K1 S2 I2 C1 F1}, diff_note={concepts: roots of unity sums | steps: expand the modulus squared; sum over roots; value; digit sum}]
Let $A(z_1)$ be a point lies on the circle whose centre is at origin and radius is $4$ unit, $z_2$ be the midpoint of $O$ and $z_1$, where $O$ is origin. Let $\omega=e^{\frac{i2\pi}{21}}$ and $\displaystyle\sum_{r=1}^{20}|z_1+\omega^rz_2|^2=K$ then the sum of digit of $K$ is divisible by...
\begin{qoptions}
\opt{A}{$3$}
\opt{B}{$4$}
\opt{C}{$5$}
\opt{D}{$7$}
\end{qoptions}
\end{bankq}
\begin{bankq}[id=TMAS04-P2-MATH-S1-Q05, source={TMAS - Test 04 (19 Nov 2023)}, paper=2, qno=5, type=multi, subtopic={Geometry}, answer={B,C,D}, answer_src=key, difficulty=8, diff_parts={K2 S2 I2 C2 F1}, diff_note={concepts: modulus of complex numbers, extrema | steps: parametrise; combine; range; check values}]
Let $z$ be complex number such that $|z|=1$ then possible values of $|1-z|+|1+z^2|$ among the following is/are
\begin{qoptions}
\opt{A}{$\dfrac{1}{\sqrt2}$}
\opt{B}{$\sqrt2$}
\opt{C}{$\pi$}
\opt{D}{$4$}
\end{qoptions}
\end{bankq}
\begin{bankq}[id=TMAS04-P2-MATH-S2-Q01, source={TMAS - Test 04 (19 Nov 2023)}, paper=2, qno=1, type=single, subtopic={Roots of unity}, answer={2}, answer_src=key, difficulty=5, diff_parts={K1 S2 I2 C1 F0}, diff_note={concepts: roots of unity | steps: expand each modulus; use conjugates; minimise; distinct naturals}]
If $a,b,c$ be three distinct natural numbers. Then the minimum value of
$\left|a+be^{i\frac{2\pi}{3}}+ce^{i\frac{4\pi}{3}}\right|^2+\left|a+be^{-i\frac{2\pi}{3}}+ce^{-i\frac{4\pi}{3}}\right|^2$ is
\end{bankq}
\begin{bankq}[id=TMAS04-P2-MATH-S2-Q04, source={TMAS - Test 04 (19 Nov 2023)}, paper=2, qno=4, type=single, subtopic={Roots of unity}, answer={1}, answer_src=key, difficulty=5, diff_parts={K1 S2 I2 C1 F0}, diff_note={concepts: polygon from roots of unity | steps: polygon of twelve vertices; area formula; value}]
If $\alpha_1,\alpha_2,\alpha_3,\ldots,\alpha_{11}$ are roots of the equation $x^{11}+x^{10}+\ldots\ldots+x+1=0$ and the area of the polygon whose vertices are $1,\alpha_1,\alpha_2,\alpha_3,\ldots,\alpha_{11}$ is $\lambda$ square unit. Then the value of $\lambda$ is
\end{bankq}

\section*{Star Review (Other) papers}
\begin{bankq}[id=STAR03-P1-MATH-S1-Q02, source={Star Review (Other) - Test 03 (04 Dec 2023)}, paper=1, qno=2, type=single, subtopic={Geometry}, answer={D}, answer_src=key, difficulty=7, diff_parts={K1 S3 I2 C2 F0}, diff_note={concepts: locus from an argument condition | steps: real and imaginary parts; slope condition; elimination; conic; latus rectum}]
If locus of point $P(z)$ which satisfies $\arg(3+z+\bar z+|z|i)=\dfrac\pi4$ is a part of conic $C$, then length of latus rectum of '$C$', is
\begin{qoptions}
\opt{A}{$1$}
\opt{B}{$\sqrt3$}
\opt{C}{$2\sqrt3$}
\opt{D}{$6$}
\end{qoptions}
\end{bankq}
\begin{bankpassage}[id=STAR03-P1-MATH-S3-G3-4, source={Star Review (Other) - Test 03 (04 Dec 2023)}, questions={3-4}]
Let $C_1$ and $C_2$ be the two curves on the complex plane defined as\par
$C_1:z+\bar z=2\,|z-1|$\par
$C_2:\arg(z+1+i)=\alpha$\par
where $\alpha$ belongs to the interval $(0,\pi)$ such that curves $C_1$ and $C_2$ touches each other at $P(z_0)$.
\end{bankpassage}
\begin{bankq}[id=STAR03-P1-MATH-S3-Q03, source={Star Review (Other) - Test 03 (04 Dec 2023)}, paper=1, qno=3, type=numerical, subtopic={Geometry}, answer={7.00}, answer_src=key, difficulty=8, diff_parts={K1 S3 I2 C2 F1}, diff_note={concepts: tangency of two curves on the complex plane | steps: identify the parabola; the ray tangent; point of contact; move to Q; sum of coordinates}, group=STAR03-P1-MATH-S3-G3-4]
A particle starts from a point $P(z_0)$. It moves horizontally away from origin by $2$ units and then vertically away from origin by $3$ units to reach at a point $Q(z_1)$. If $z_1=x_1+iy_1$, then $(x_1+y_1)$ equals -
\end{bankq}
\begin{bankq}[id=STAR03-P1-MATH-S3-Q04, source={Star Review (Other) - Test 03 (04 Dec 2023)}, paper=1, qno=4, type=numerical, subtopic={Geometry}, answer={0.665}, answer_src=key, difficulty=7, diff_parts={K1 S2 I2 C2 F1}, diff_note={concepts: area bounded by a parabola and a chord after rotation | steps: point of contact; rotated point; chord; area with the parabola}, group=STAR03-P1-MATH-S3-G3-4]
If $P(z_0)$ is rotated about origin through an angle '$2\alpha$' in clockwise direction, then the area bounded (in sq. unit) by $C_1$ and the line joining $P(z_0)$ and $Q(z_0')$ is (where $Q(z_0')$ is the new position $P(z_0)$ after rotation) sq. unit-
\end{bankq}

\section*{Nurture (General) papers}
\begin{bankq}[id=NUR04-P1-MATH-S2-Q09, source={Nurture (General) - Test 04 (11 Feb 2024)}, paper=1, qno=9, type=multi, subtopic={Geometry}, answer={A,D}, answer_src=solved, working={x/a+y/b=1 with x=(z+zbar)/2, y=(z-zbar)/2i}, difficulty=5, diff_parts={K1 S2 I1 C1 F1}, diff_note={concepts: complex form of a line | steps: two intercept form; substitute; compare}]
Equation of line through a and ib such that $a,b\in R$ and $a,b\ne0$ is
\begin{qoptions}
\opt{A}{$z\left(\dfrac1{2a}+\dfrac1{2ib}\right)+\overline z\left(\dfrac1{2a}-\dfrac1{2ib}\right)=1$}
\opt{B}{$z\left(\dfrac1{2a}+\dfrac i{2b}\right)+\overline z\left(\dfrac1{2a}+\dfrac i{2b}\right)=1$}
\opt{C}{$z\left(\dfrac1{2a}-\dfrac i{2b}\right)+\overline z\left(\dfrac1{2a}-\dfrac i{2b}\right)=1$}
\opt{D}{$z\left(\dfrac1{2a}-\dfrac i{2b}\right)+\overline z\left(\dfrac1{2a}+\dfrac i{2b}\right)=1$}
\end{qoptions}
\end{bankq}
\begin{bankq}[id=NUR04-P1-MATH-S3-Q15, source={Nurture (General) - Test 04 (11 Feb 2024)}, paper=1, qno=15, type=numerical, subtopic={Geometry}, answer={2}, answer_src=solved, working={|PC2|-|PC1|=2 defines a hyperbola when |z1-z2|>2}, difficulty=5, diff_parts={K1 S2 I1 C1 F1}, diff_note={concepts: locus of centres of touching circles | steps: distance conditions; difference; hyperbola condition}]
The locus of the centre of a circle which touches the circles $|z-z_1|=2$ and $|z-z_2|=4$, externally is a hyperbola if $|z_1-z_2|$, is greater than \underline{\hspace{1cm}}
\end{bankq}
\end{document}
